% tableStyle.tex — Single source of truth for table/figure colors and
% header/zebra styling used across the whole thesis.
%
% Palette: Okabe & Ito (2008), "Color Universal Design" — the standard
% categorical palette distinguishable under protanopia, deuteranopia and
% tritanopia. Every table/gantt-chart color in the thesis must come from
% here; do not \definecolor ad hoc inside chapters/tables/*.tex.

% ── Okabe-Ito base palette ────────────────────────────────────────────────
\definecolor{cbOrange}{RGB}{230,159,0}
\definecolor{cbSkyBlue}{RGB}{86,180,233}
\definecolor{cbGreen}{RGB}{0,158,115}     % bluish green
\definecolor{cbYellow}{RGB}{240,228,66}
\definecolor{cbBlue}{RGB}{0,114,178}
\definecolor{cbVermillion}{RGB}{213,94,0}
\definecolor{cbPurple}{RGB}{204,121,167}
\definecolor{cbGrey}{RGB}{166,166,166}

% ── Semantic roles (used by every table in the thesis) ───────────────────
% Header row: solid blue background, white bold text.
\colorlet{tableHeaderColor}{cbBlue}
% Zebra striping on long/data-heavy tables: a faint tint of the header hue,
% so it reads purely as a lightness cue (works for every CVD type) while
% still tying visually back to the header color.
\colorlet{tableAltRowColor}{cbBlue!6}
% Wavelet coefficient examples (tables/{regular,packet}WaveletExample.tex):
% grey = raw signal label, sky blue = row/column index label,
% blue = approximation band (cA), orange = detail band (cD).
% Blue/orange is the single most robust 2-hue pair across all CVD types.
\colorlet{tableSignalLabelColor}{cbGrey}
\colorlet{tableIndexLabelColor}{cbSkyBlue}
\colorlet{tableSignalValuesColor}{cbGreen!35}
\colorlet{tableApproxBandColor}{cbBlue!35}
\colorlet{tableDetailBandColor}{cbOrange!35}

% Gantt chart (chapters/08-proposedApproach.tex, tab:cronogramaGantt):
% same palette, kept here so the chart matches the tables it sits beside.
\colorlet{ganttBarColor}{cbGreen}
\colorlet{ganttMilestoneColor}{cbGreen}
\colorlet{ganttTitleColor}{cbBlue!30}
\colorlet{ganttHighlightColor}{cbOrange}

% ── Reusable header row ────────────────────────────────────────────────────
% Usage inside a tabular/longtable, as the first thing in a header row:
%   \tableHeaderRow \textbf{A} & \textbf{B} \\
\newcommand{\tableHeaderRow}{\rowcolor{tableHeaderColor}\color{white}}

% ── Zebra striping for table/longtable bodies ─────────────────────────────
% Call once, right before \begin{tabular}/\begin{longtable}. Colors rows
% 2, 4, 6, ... (i.e. every data row after a 1-row header) with
% tableAltRowColor; row 1 (the header) is left untouched so \tableHeaderRow
% still controls it. NOTE: \rowcolor itself must always be the literal first
% token of a row — never wrap it in another macro (e.g. a counter-stepper),
% since colortbl's \noalign trick requires nothing to precede it.
\newcommand{\tableZebraStripe}{\rowcolors{2}{white}{tableAltRowColor}}
